% =====================================================================
%  Chương 6 — Kết luận
% =====================================================================
\chapter{Kết luận}
\label{ch:ketluan}

\section{Trả lời năm mục của đề}

\begin{enumerate}
  \item \textbf{Khái niệm cơ bản qua hàm và hợp hàm} (Chương~\ref{ch:khainiem}). Neuron, hàm kích hoạt, tích chập, pooling,
        softmax và cross-entropy được viết thành các hàm NumPy trong \texttt{a05/concepts.py}. CNN là hợp của các hàm đó, và
        lan truyền ngược là quy tắc đạo hàm của hàm hợp. \SL{concepts/checks_passed}/\SL{concepts/checks_total} phép kiểm chứng
        khớp với PyTorch; gradient viết tay của tầng tích chập sai lệch so với sai phân hữu hạn chưa tới $10^{-6}$.
  \item \textbf{Các mô hình phát triển} (Chương~\ref{ch:kientruc}). Chín mốc từ LeNet-5 đến ViT, mỗi mốc là một cơ chế viết
        được dưới dạng hàm: chồng conv $3\times3$, nhánh song song, $F(x)+x$, nối mọi tầng, tách depthwise, nhân trọng số kênh,
        self-attention. Mỗi cơ chế có một lớp PyTorch trong \texttt{a05/blocks.py}, đã kiểm tra shape và số tham số.
  \item \textbf{Dữ liệu} (Chương~\ref{ch:dulieu}). CIFAR-10, CIFAR-100 (bài hiểu ``MNIST-100'' của đề là tập này) và Diabetes
        Prediction, kèm đường dẫn đã kiểm tra, thống kê mô tả, cách chia tập và dạng tensor. Diabetes bị bỏ
        \SL{data/diabetes/n_dup} dòng trùng trước khi chia, và mọi phép biến đổi chỉ học trên nhánh train.
  \item \textbf{Một CNN cơ bản và ba mô hình phát triển} (Chương~\ref{ch:mohinh}). BasicCNN, VGGNet, ResNet, SE-ResNet, mỗi
        bước thêm đúng một cơ chế, có bản 2D cho ảnh và bản 1D cho Diabetes từ cùng một định nghĩa. Mười hai lượt huấn luyện
        (3 tập $\times$ 4 mô hình) dùng chung một cấu hình trong mỗi tập.
  \item \textbf{So sánh và đánh giá} (Chương~\ref{ch:sosanh}). So theo chất lượng, chi phí, hội tụ, lỗi và biểu diễn. Kết quả
        tốt nhất: \SL{cifar10/models/resnet/acc}\,\% trên CIFAR-10 (ResNet), \SL{cifar100/models/seresnet/acc}\,\% top-1 và
        \SL{cifar100/models/seresnet/top5}\,\% top-5 trên CIFAR-100 (SE-ResNet), ROC-AUC \SL{diabetes/models/basic/roc_auc}
        trên Diabetes (BasicCNN).
\end{enumerate}

\section{Điều rút ra}

Kết quả quan trọng nhất của bài không phải mô hình nào đứng đầu, mà là \emph{một cơ chế chỉ có ích khi mô hình đang gặp
đúng vấn đề mà cơ chế đó sinh ra để giải quyết}. BatchNorm được tạo ra để huấn luyện ổn định, và M0, mô hình duy nhất không có
nó, là mô hình đã phân kỳ ở tốc độ học cao khi chưa cắt gradient. Đường tắt giải quyết bài toán tối ưu mạng sâu, nên nó không đổi gì khi M1 đã gần như
thuộc lòng CIFAR-10, và chỉ thêm được \SL{compare/delta/cifar100/vgg_resnet} điểm trên CIFAR-100. Cả ba cơ chế đều khai thác cấu trúc không gian của ảnh,
nên không giúp gì cho một hàng dữ liệu bảng không có cấu trúc đó. Đọc lịch sử CNN như một chuỗi ``kiến trúc sau luôn tốt hơn
kiến trúc trước'' là đọc sai: mỗi kiến trúc tốt hơn \emph{trong điều kiện} mà nó được thiết kế.

\section{Hạn chế}

\begin{itemize}
  \item \textbf{Một hạt giống.} Mỗi mô hình chạy một lần. Chênh lệch dưới 0,5 điểm không được dùng để kết luận, và ngay cả
        chênh lệch \SL{compare/delta/cifar100/vgg_resnet} điểm của đường tắt trên CIFAR-100 cũng chỉ nên coi là dấu hiệu.
  \item \textbf{Bước M0$\to$M1 gộp ba thay đổi.} Thí nghiệm không tách được phần đóng góp riêng của độ sâu, BatchNorm và GAP.
  \item \textbf{Mạng nông.} Với 12 tầng tích chập, thí nghiệm không kiểm tra được lợi ích của đường tắt ở độ sâu mà nó được
        thiết kế cho.
  \item \textbf{Thời gian đo trên GPU laptop bị hạ xung vì nhiệt.} Số giây mỗi epoch chỉ so sánh được giữa các mô hình với nhau.
  \item \textbf{``MNIST-100''} được hiểu là CIFAR-100. Nếu đề muốn một tập khác, Chương~\ref{ch:dulieu} và phần CIFAR-100 của
        Chương~\ref{ch:sosanh} cần làm lại với tập đó; mã nguồn không phải đổi gì ngoài hàm nạp dữ liệu.
\end{itemize}
